% =========================================================
% CHAPITRE 5 : SPRINT 3
% =========================================================

\clearpage

\refstepcounter{chapter}

\setcounter{section}{0}
\setcounter{subsection}{0}
\setcounter{subsubsection}{0}

\phantomsection
\addcontentsline{toc}{chapter}
{Chapitre \thechapter : Sprint 3 -- Gestion du cycle de vie des modèles}

\markboth{Chapitre \thechapter}{Chapitre \thechapter}

\thispagestyle{empty}


% =========================================================
% PAGE D'OUVERTURE + PLAN
% =========================================================

\vspace*{0.8cm}

\begin{center}

\rule{0.30\textwidth}{1.2pt}
\hspace{0.5cm}
\textit{Chapitre \thechapter}
\hspace{0.5cm}
\rule{0.30\textwidth}{1.2pt}

\vspace{0.9cm}

\rule{\textwidth}{0.8pt}

\vspace{0.35cm}

{\fontsize{21}{24}\selectfont
\textsc{
\parbox{0.95\textwidth}{
\centering
Chapitre \thechapter : Sprint 3 -- Gestion du cycle de vie\\
des modèles
}
}
}

\vspace{0.35cm}

\rule{\textwidth}{0.8pt}

\end{center}


\vspace{1.4cm}

{\Large\bfseries Plan}

\vspace{0.65cm}

\begin{center}

\renewcommand{\arraystretch}{1.35}

\begin{tabular}{p{0.8cm}p{10.7cm}r}

\textbf{1}
&
\textbf{Introduction}
&
\pageref{sec:ch5-introduction}
\\[0.22cm]

\textbf{2}
&
\textbf{Objectifs et travaux du Sprint 3}
&
\pageref{sec:sprint3-objectifs}
\\[0.22cm]

\textbf{3}
&
\textbf{Versionnement des modèles}
&
\pageref{sec:sprint3-versioning}
\\[0.22cm]

\textbf{4}
&
\textbf{Déploiement des modèles}
&
\pageref{sec:sprint3-deployment}
\\[0.22cm]

\textbf{5}
&
\textbf{Services de prédiction}
&
\pageref{sec:sprint3-prediction}
\\[0.22cm]

\textbf{6}
&
\textbf{Gestion des accès}
&
\pageref{sec:sprint3-access}
\\[0.22cm]

\textbf{7}
&
\textbf{Monitoring}
&
\pageref{sec:sprint3-monitoring}
\\[0.22cm]

\textbf{8}
&
\textbf{Stratégies de réentraînement}
&
\pageref{sec:sprint3-retraining}
\\[0.22cm]

\textbf{9}
&
\textbf{Intégration continue avec GitHub Actions}
&
\pageref{sec:sprint3-cicd}
\\[0.22cm]

\textbf{10}
&
\textbf{Bilan du Sprint 3}
&
\pageref{sec:sprint3-bilan}
\\[0.22cm]

\textbf{11}
&
\textbf{Conclusion}
&
\pageref{sec:ch5-conclusion}
\\[0.22cm]

\end{tabular}

\renewcommand{\arraystretch}{1}

\end{center}

\clearpage


% =========================================================
% 5.1 INTRODUCTION
% =========================================================

\section{Introduction}
\label{sec:ch5-introduction}

Après avoir construit le moteur d'entraînement et le parcours
utilisateur lors des sprints précédents, le troisième sprint s'est
concentré sur le \textbf{cycle de vie complet des modèles}.

L'objectif était de permettre à un modèle d'aller de l'entraînement
jusqu'à la production en passant par le versionnement, le déploiement,
la prédiction et le monitoring, tout en disposant de mécanismes
automatisés de réentraînement.

Ce sprint couvre donc l'ensemble des étapes qui transforment un run
d'entraînement en un service de prédiction opérationnel, surveillé et
maintenu dans le temps.


% =========================================================
% 5.2 OBJECTIFS ET TRAVAUX
% =========================================================

\section{Objectifs et travaux du Sprint 3}
\label{sec:sprint3-objectifs}

Le Sprint 3 s'est déroulé du \textbf{18 mai au 30 juin}. Les principaux
travaux réalisés pendant cette période sont résumés dans le tableau
\ref{tab:travaux-sprint3}.

\begin{table}[H]

\centering
\small

\renewcommand{\arraystretch}{1.16}
\setlength{\tabcolsep}{5pt}

\begin{tabular}{|p{4cm}|p{9.2cm}|}

\hline

\rowcolor{cyan!20}
\textbf{Élément}
&
\textbf{Rôle dans le Sprint 3}
\\
\hline

\textbf{Versionnement}
&
Gérer les versions des modèles via MLflow Model Registry, promouvoir
la meilleure version en production grâce au score global.
\\
\hline

\rowcolor{cyan!7}
\textbf{Déploiement}
&
Déployer les modèles sur Kubernetes avec une stratégie Blue/Green
pour un basculement sans interruption de service.
\\
\hline

\textbf{Services de prédiction}
&
Exposer des endpoints de prédiction internes (JWT) et externes
(API key) pour tabulaire, texte et images.
\\
\hline

\rowcolor{cyan!7}
\textbf{Gestion des accès}
&
Contrôler la visibilité des déploiements (privé, partagé, public)
et gérer les rôles utilisateur et administrateur.
\\
\hline

\textbf{Monitoring}
&
Suivre l'état des modèles déployés, détecter les dérives de données,
les baisses de performance et les erreurs de prédiction.
\\
\hline

\rowcolor{cyan!7}
\textbf{Réentraînement}
&
Relancer automatiquement l'entraînement en cas de dérive ou de
baisse de qualité, ou suite à un push GitHub.
\\
\hline

\textbf{CI/CD}
&
Automatiser la construction et le push des images Docker via
GitHub Actions.
\\
\hline

\end{tabular}

\caption{Principaux travaux réalisés lors du Sprint 3}
\label{tab:travaux-sprint3}

\end{table}


% =========================================================
% 5.3 VERSIONNEMENT DES MODÈLES
% =========================================================

\section{Versionnement des modèles}
\label{sec:sprint3-versioning}

Le versionnement permet de conserver une trace de chaque modèle
entraîné et de comparer les versions entre elles. Il repose sur le
\textbf{MLflow Model Registry} complété par une synchronisation
vers la base de données MySQL de la plateforme.


% =========================================================
\subsection{MLflow Model Registry}
\label{subsec:sprint3-mlflow-registry}

Chaque run terminé enregistre automatiquement le modèle dans le
\textit{MLflow Model Registry} sous un nom de groupe correspondant
au projet. Chaque enregistrement reçoit un numéro de version incrémental.

Deux alias sont utilisés pour identifier les versions importantes :
\texttt{@latest} pour la version la plus récente et
\texttt{@production} pour la version déployée en production.

Le service \texttt{MlflowRegistryClientService} interroge l'API
REST de MLflow pour résoudre la version en production. Il utilise
d'abord l'endpoint d'alias, puis effectue une recherche de fallback
si l'alias n'est pas trouvé.


% =========================================================
\subsection{Synchronisation MLflow -- MySQL}
\label{subsec:sprint3-sync}

Afin de ne pas dépendre uniquement de MLflow pour l'affichage des
informations, le backend synchronise régulièrement les métadonnées
du Registry vers MySQL.

Cette synchronisation est déclenchée par le service
\texttt{ModelRegistryService} et met à jour les colonnes
\texttt{modelVersion} et \texttt{productionVersion} de la table
\texttt{MlExperiment}.

\begin{figure}[H]
    \centering

    \includegraphics[
        width=0.85\textwidth,
        keepaspectratio
    ]{images/sprint3_model_registry.png}

    \caption{Interface de gestion des modèles enregistrés}
    \label{fig:sprint3-model-registry}
\end{figure}

La figure \ref{fig:sprint3-model-registry} présente la liste des
modèles enregistrés avec leur version actuelle, la version en
production et le score global associé.


% =========================================================
\subsection{Promotion automatique}
\label{subsec:sprint3-promotion}

Lorsqu'un run se termine avec succès, le service
\texttt{ModelPromotionService} évalue si la nouvelle version doit
être promue en production.

La décision repose sur le \textbf{score global} calculé à partir des
métriques prioritaires définies par l'utilisateur. Si la nouvelle
version obtient un score supérieur à la version actuelle, elle est
automatiquement promue en production via l'alias \texttt{@production}
du MLflow Model Registry.

Cette promotion est également synchronisée vers MySQL pour garantir
la cohérence entre les deux sources de données.


% =========================================================
% 5.4 DÉPLOIEMENT DES MODÈLES
% =========================================================

\section{Déploiement des modèles}
\label{sec:sprint3-deployment}

Le déploiement transforme un modèle enregistré en un service de
prédiction accessible via Kubernetes. La plateforme utilise une
stratégie \textbf{Blue/Green} pour garantir un basculement sans
interruption de service.


% =========================================================
\subsection{Stratégie Blue/Green}
\label{subsec:sprint3-blue-green}

La stratégie Blue/Green est une technique de déploiement qui consiste
à maintenir \textbf{deux environnements identiques} en permanence :
un environnement \textbf{blue} (actif) qui reçoit tout le traffic
et un environnement \textbf{green} (standby) qui reste prêt à
prendre le relais.

Contrairement au déploiement classique où l'ancienne version est
arrêtée avant le démarrage de la nouvelle (causant une interruption),
la stratégie Blue/Green permet de déployer la nouvelle version
\textbf{sans aucune coupure de service}. Le traffic n'est basculé
vers le nouveau track qu'après validation complète.

Dans notre plateforme, chaque track correspond à un ensemble de
ressources Kubernetes :

\begin{itemize}
    \item un \textbf{Deployment} contenant le conteneur de serving
    du modèle ;
    \item un \textbf{Service} ClusterIP dont le sélecteur pointe
    vers le track actif via l'étiquette \texttt{serving-track}.
\end{itemize}

À tout moment, le Service ne route le traffic que vers un seul
track. L'autre reste disponible mais ne reçoit aucune requête.

\begin{figure}[H]
    \centering

    \includegraphics[
        width=0.85\textwidth,
        keepaspectratio
    ]{images/sprint3_k8s_deployments.png}

    \caption{Déploiements et services des modèles dans le cluster Kubernetes}
    \label{fig:sprint3-k8s-deployments}
\end{figure}

La figure \ref{fig:sprint3-k8s-deployments} montre les déploiements
et services des modèles dans le cluster Kubernetes. Chaque modèle
dispose d'un Deployment (par exemple \texttt{class-model-v3-deployment})
et d'un Service stable (par exemple \texttt{class-model-service}) qui
route le traffic vers le track actif.


% =========================================================
\subsection{Processus de déploiement}
\label{subsec:sprint3-deploy-process}

Le déploiement d'un nouveau modèle suit les étapes suivantes :

\begin{enumerate}
    \item \textbf{Résolution du modèle} : le backend identifie le
    nom du modèle, la version à déployer et l'image Docker de
    serving correspondante (image standard ou image personnalisée
    construite via Kaniko).

    \item \textbf{Déploiement sur le track inactif} : le service
    \texttt{KubernetesModelServingService} crée un nouveau
    Deployment sur le track qui n'est pas actif (par exemple
    \texttt{green} si \texttt{blue} est en production). Le
    Deployment reçoit la variable d'environnement
    \texttt{MODEL\_URI=models:/nom-modele/version} qui indique
    au conteneur de serving quel modèle charger depuis MLflow.

    \item \textbf{Attente de la disponibilité} : le backend
    attend que les pods du nouveau Deployment soient prêts
    (\texttt{isModelServerReady()}), c'est-à-dire que les
    conteneurs soient démarrés et passent les probes de
    santé Kubernetes.

    \item \textbf{Smoke test} : avant de basculer le traffic,
    un test de validation est envoyé directement au endpoint
    de prédiction du nouveau track. Ce test consiste à envoyer
    une requête de prédiction factice et à vérifier que le
    modèle répond correctement avec un code HTTP 200 et une
    réponse valide.

    \item \textbf{Basculement du traffic} : si le smoke test
    est concluant, le Service Kubernetes est mis à jour pour
    pointer vers le nouveau track (\texttt{serving-track=green}).
    Le traffic est instantanément routé vers la nouvelle version
    sans aucune coupure.

    \item \textbf{Nettoyage} : l'ancien Deployment (track
    précédent) est conservé temporairement en standby puis
    supprimé après confirmation que le nouveau fonctionne
    correctement.
\end{enumerate}


% =========================================================
\subsection{Le smoke test}
\label{subsec:sprint3-smoke-test}

Le \textbf{smoke test} (littéralement « test de fumée ») est une
vérification minimale qui consiste à s'assurer qu'un service est
opérationnel avant de lui confier du traffic réel.

Dans le contexte du déploiement de modèles, le smoke test vérifie
que :

\begin{itemize}
    \item le conteneur de serving est démarré et répond aux requêtes ;
    \item le modèle a été correctement chargé depuis MLflow via
    la variable \texttt{MODEL\_URI} ;
    \item une prédiction peut être effectuée sans erreur ;
    \item le temps de réponse est dans des limites acceptables.
\end{itemize}

Ce test est déclenché par la méthode
\texttt{probePredictionEndpoint()} du service
\texttt{KubernetesModelServingService}. Si le test échoue, le
déploiement est annulé et l'ancien track reste actif, garantissant
ainsi la continuité de service.

Le smoke test est un garde-fou essentiel : il évite de basculer
le traffic vers un modèle défaillant ou un conteneur qui n'a pas
réussi à se démarrer correctement.


% =========================================================
\subsection{Auto-déploiement}
\label{subsec:sprint3-auto-deploy}

Le service \texttt{AutoProductionDeploymentReconcilerService} surveille
les changements d'état de production. Lorsqu'une nouvelle version est
promue en production (via le score global), il déclenche
automatiquement le déploiement sans intervention manuelle.

Ce mécanisme assure que la plateforme reste toujours synchronisée
avec la version recommandée du modèle. L'utilisateur peut également
déclencher manuellement un déploiement depuis l'interface si nécessaire.


% =========================================================
% 5.5 SERVICES DE PRÉDICTION
% =========================================================

\section{Services de prédiction}
\label{sec:sprint3-prediction}

La plateforme expose deux voies de prédiction : une voie interne
authentifiée par JWT et une voie externe sécurisée par une API key.


% =========================================================
\subsection{Prédiction interne}
\label{subsec:sprint3-prediction-internal}

La prédiction interne est accessible via les endpoints
\texttt{/api/predictions/predict},
\texttt{/api/predictions/predict-image} et
\texttt{/api/predictions/predict-text}.

L'utilisateur sélectionne un modèle déployé parmi ceux auxquels il a
accès, soumet ses données et reçoit la prédiction accompagnée du
nom du modèle, de la version et du latence.

Chaque prédiction est enregistrée dans la table
\texttt{PredictionLog} pour alimenter le monitoring.

\begin{figure}[H]
    \centering

    \includegraphics[
        width=0.85\textwidth,
        keepaspectratio
    ]{images/sprint3_prediction_internal.png}

    \caption{Interface de prédiction interne}
    \label{fig:sprint3-prediction-internal}
\end{figure}


% =========================================================
\subsection{SMART Model Gateway (prédiction externe)}
\label{subsec:sprint3-smart-gateway}

Le \textbf{SMART Model Gateway} est un endpoint externe dédié à
l'intégration avec des applications tierces. Il est accessible via
\texttt{/api/external/models/\{nom\}/predict} et
\texttt{/api/external/models/\{nom\}/predict-image}.

L'accès est sécurisé par une \textbf{API key} générée par le
propriétaire du modèle. Le gateway valide la clé, résout le
déploiement actif, appelle le modèle et retourne la réponse
accompagnée des métadonnées (version, ID de déploiement, ID de
requête).

Un compteur atomique suit l'utilisation de chaque API key pour
facturation ou monitoring.

\begin{figure}[H]
    \centering

    \includegraphics[
        width=0.85\textwidth,
        keepaspectratio
    ]{images/sprint3_smart_gateway.png}

    \caption{SMART Model Gateway -- endpoint externe de prédiction}
    \label{fig:sprint3-smart-gateway}
\end{figure}


% =========================================================
% 5.6 GESTION DES ACCÈS
% =========================================================

\section{Gestion des accès}
\label{sec:sprint3-access}

La gestion des accès repose sur deux niveaux : les rôles utilisateur
et la visibilité des déploiements.


% =========================================================
\subsection{Rôles et permissions}
\label{subsec:sprint3-roles}

Deux rôles principaux ont été définis :

\begin{itemize}
    \item l'\textbf{utilisateur standard} peut créer des projets,
    entraîner des modèles, les déployer et les utiliser en prédiction ;
    \item l'\textbf{administrateur} peut en outre gérer les comptes,
    consulter la vue globale du monitoring et recevoir les alertes
    critiques.
\end{itemize}

Côté frontend, les guards Angular (\texttt{authGuard},
\texttt{adminGuard}) contrôlent l'accès aux pages. Côté backend,
les endpoints sensibles vérifient le rôle de l'utilisateur via le
token JWT.


% =========================================================
\subsection{Visibilité des déploiements}
\label{subsec:sprint3-visibility}

Chaque déploiement dispose d'une visibilité qui contrôle qui peut
l'utiliser pour des prédictions :

\begin{table}[H]

\centering
\small

\renewcommand{\arraystretch}{1.15}

\begin{tabular}{|p{3cm}|p{10cm}|}

\hline

\rowcolor{cyan!20}
\textbf{Visibilité}
&
\textbf{Règle d'accès}
\\
\hline

\textbf{PRIVATE}
&
Seul le propriétaire du projet peut accéder au déploiement.
\\
\hline

\rowcolor{cyan!7}
\textbf{SHARED}
&
Le propriétaire et les utilisateurs spécifiquement partagés peuvent
accéder.
\\
\hline

\textbf{PUBLIC API}
&
Tout utilisateur authentifié peut accéder au déploiement via le
SMART Gateway.
\\
\hline

\end{tabular}

\caption{Visibilité des déploiements et règles d'accès}
\label{tab:sprint3-visibility}

\end{table}

Le service \texttt{ModelAccessControlService} centralise la logique
de vérification des accès pour les opérations de prédiction, de
gestion des clés API et de partage.


% =========================================================
% 5.7 MONITORING
% =========================================================

\section{Monitoring}
\label{sec:sprint3-monitoring}

Le monitoring assure la surveillance continue de l'état des modèles
déployés, de la qualité des prédictions et de la santé du cluster
Kubernetes.


% =========================================================
\subsection{Architecture de monitoring}
\label{subsec:sprint3-monitoring-arch}

L'architecture de monitoring repose sur trois piliers :

\begin{itemize}
    \item \textbf{Prometheus} collecte les métriques exposées par le
    backend via l'endpoint \texttt{/actuator/prometheus} et les métriques
    Kubernetes ;
    \item \textbf{Grafana} fournit les tableaux de bord pour la
    visualisation ;
    \item \textbf{Alertmanager} gère les alertes et les notifications.
\end{itemize}

Un \texttt{ServiceMonitor} Kubernetes configure la collecte automatique
toutes les 15 secondes. Les métriques métier sont enregistrées via
Micrometer dans le service
\texttt{MlopsBusinessMetricsService}.

\begin{figure}[H]
    \centering

    \includegraphics[
        width=0.85\textwidth,
        keepaspectratio
    ]{images/sprint3_monitoring_architecture.png}

    \caption{Architecture du monitoring Prometheus -- Grafana}
    \label{fig:sprint3-monitoring-architecture}
\end{figure}


% =========================================================
\subsection{Métriques métier}
\label{subsec:sprint3-business-metrics}

Le backend expose plusieurs métriques via Micrometer :

\begin{itemize}
    \item \texttt{mlops\_experiments\_total} -- nombre total
    d'expériences ;
    \item \texttt{mlops\_experiments\_running} -- expériences en cours ;
    \item \texttt{mlops\_deployments\_deployed} -- déploiements actifs ;
    \item \texttt{mlops\_active\_production\_models} -- modèles en
    production ;
    \item \texttt{mlops\_public\_deployments} -- déploiements publics.
\end{itemize}

Ces métriques sont collectées par Prometheus et affichées dans Grafana.

\begin{figure}[H]
    \centering

    \includegraphics[
        width=0.85\textwidth,
        keepaspectratio
    ]{images/sprint3_grafana_dashboard.png}

    \caption{Tableau de bord Grafana -- métriques de la plateforme}
    \label{fig:sprint3-grafana-dashboard}
\end{figure}


% =========================================================
\subsection{Alertes et détection de dérive}
\label{subsec:sprint3-alerts}

Le service \texttt{MonitoringAlertScheduler} exécute des vérifications
périodiques pour détecter les problèmes :

\begin{itemize}
    \item échecs d'entraînement ;
    \item déploiements en échec ;
    \item latence excessive du backend ;
    \item dérive de données (\textit{data drift}) via Evidently ;
    \item taux d'erreur de prédiction élevé ;
    \item baisse de performance du modèle.
\end{itemize}

Lorsqu'un problème est détecté, le service
\texttt{MonitoringAlertService} crée une alerte idempotente et
notifie le propriétaire du projet. Les alertes critiques notifient
également les administrateurs.

\begin{figure}[H]
    \centering

    \includegraphics[
        width=0.85\textwidth,
        keepaspectratio
    ]{images/sprint3_monitoring_alerts.png}

    \caption{Alertes de monitoring et notifications}
    \label{fig:sprint3-monitoring-alerts}
\end{figure}


% =========================================================
% 5.8 STRATÉGIES DE RÉENTRAÎNEMENT
% =========================================================

\section{Stratégies de réentraînement}
\label{sec:sprint3-retraining}

La plateforme propose deux stratégies principales de réentraînement
automatisé : la stratégie \textbf{Monitoring Smart} et la stratégie
\textbf{GitHub Actions}.


% =========================================================
\subsection{Stratégie 1 : Monitoring Smart}
\label{subsec:sprint3-retraining-monitoring}

La stratégie \textbf{Monitoring Smart} déclenche un réentraînement
automatique lorsque le monitoring détecte un problème de qualité
du modèle.

Le processus suit les étapes suivantes :

\begin{enumerate}
    \item Le \texttt{MonitoringAlertScheduler} détecte une anomalie
    (dérive de données, taux d'erreur élevé, baisse de performance) ;
    \item Le \texttt{MonitoringRetrainingTriggerService} évalue si
    le réentraînement est justifié en fonction des seuils configurés ;
    \item Le \texttt{AutoRetrainingService} crée une nouvelle
    expérience de réentraînement en reprenant les données et le script
    de la version production ;
    \item L'entraînement est lancé sur Kubeflow et le nouveau modèle
    est évalué ;
    \item Si le nouveau modèle est meilleur, il est automatiquement
    promu et déployé.
\end{enumerate}

\begin{figure}[H]
    \centering

    \includegraphics[
        width=0.90\textwidth,
        keepaspectratio
    ]{images/uml_retraining_monitoring.png}

    \caption{Diagramme de séquence -- Réentraînement par monitoring}
    \label{fig:uml-retraining-monitoring}
\end{figure}

La configuration du réentraînement est stockée dans la table
\texttt{RetrainingConfig} et permet de paramétrer :

\begin{itemize}
    \item la fréquence de vérification (journalière, hebdomadaire,
    mensuelle) ;
    \item les seuils d'alerte (taux d'erreur maximal, seuil de
    performance, seuil de dérive) ;
    \item la politique de données (données originales, dernière
    version, fusion) ;
    \item le niveau de sécurité (alerter uniquement ou bloquer le
    modèle).
\end{itemize}

\begin{figure}[H]
    \centering

    \includegraphics[
        width=0.85\textwidth,
        keepaspectratio
    ]{images/sprint3_retraining_config.png}

    \caption{Configuration de la stratégie de réentraînement Monitoring Smart}
    \label{fig:sprint3-retraining-config}
\end{figure}


% =========================================================
\subsection{Stratégie 2 : GitHub Actions}
\label{subsec:sprint3-retraining-github}

La stratégie \textbf{GitHub Actions} déclenche un réentraînement
lorsqu'un push est effectué sur le dépôt GitHub associé au projet.

Le processus suit les étapes suivantes :

\begin{enumerate}
    \item L'utilisateur configure le webhook GitHub dans la plateforme,
    en spécifiant le dépôt, la branche et en générant un secret HMAC ;
    \item Lors d'un push sur la branche configurée, GitHub envoie
    un webhook vers le backend ;
    \item Le \texttt{GitHubRetrainingService} valide la signature
    HMAC et extrait le commit SHA ;
    \item Le \texttt{AutoRetrainingService} lance un réentraînement
    avec les données à jour du dépôt ;
    \item Le modèle résultant est évalué et promu si nécessaire.
\end{enumerate}

\begin{figure}[H]
    \centering

    \includegraphics[
        width=0.90\textwidth,
        keepaspectratio
    ]{images/uml_retraining_github.png}

    \caption{Diagramme de séquence -- Réentraînement via GitHub Actions}
    \label{fig:uml-retraining-github}
\end{figure}

\begin{figure}[H]
    \centering

    \includegraphics[
        width=0.85\textwidth,
        keepaspectratio
    ]{images/sprint3_github_retraining_config.png}

    \caption{Configuration du réentraînement GitHub}
    \label{fig:sprint3-github-retraining-config}
\end{figure}


% =========================================================
\subsection{Comparaison des deux stratégies}
\label{subsec:sprint3-retraining-comparison}

\begin{table}[H]

\centering
\small

\renewcommand{\arraystretch}{1.15}
\setlength{\tabcolsep}{4pt}

\begin{tabular}{|p{3.5cm}|p{4.8cm}|p{4.8cm}|}

\hline

\rowcolor{cyan!20}
\textbf{Critère}
&
\textbf{Monitoring Smart}
&
\textbf{GitHub Actions}
\\
\hline

\textbf{Déclencheur}
&
Dérive de données ou baisse de performance
&
Push sur le dépôt GitHub
\\
\hline

\rowcolor{cyan!7}
\textbf{Automatique}
&
Oui, basé sur les seuils
&
Oui, basé sur les commits
\\
\hline

\textbf{Données utilisées}
&
Données de production récentes
&
Dernière version du dépôt
\\
\hline

\rowcolor{cyan!7}
\textbf{Cas d'usage}
&
Maintenance continue du modèle
&
Évolution du code ou des données
\\
\hline

\textbf{Configuration}
&
Seuils et fréquence dans la plateforme
&
Dépôt, branche et webhook GitHub
\\
\hline

\end{tabular}

\caption{Comparaison des stratégies de réentraînement}
\label{tab:sprint3-retraining-comparison}

\end{table}


% =========================================================
% 5.9 INTÉGRATION CONTINUE AVEC GITHUB ACTIONS
% =========================================================

\section{Intégration continue avec GitHub Actions}
\label{sec:sprint3-cicd}

L'intégration continue automatisé la construction et le déploiement
des images Docker via un workflow GitHub Actions.


% =========================================================
\subsection{Workflow de build}
\label{subsec:sprint3-workflow}

Le workflow \texttt{docker-build.yml} se déclenche à chaque push sur
la branche \texttt{main}. Il construit deux images Docker :

\begin{itemize}
    \item \textbf{Backend Spring Boot} : compilation avec Java 17,
    build de l'image Docker, push vers Docker Hub ;
    \item \textbf{Frontend Angular} : build avec Node.js 20,
    construction de l'image Nginx, push vers Docker Hub.
\end{itemize}

\begin{figure}[H]
    \centering

    \includegraphics[
        width=0.90\textwidth,
        keepaspectratio
    ]{images/sprint3_github_actions_workflow.png}

    \caption{Workflow GitHub Actions -- build et push des images Docker}
    \label{fig:sprint3-github-actions-workflow}
\end{figure}


% =========================================================
\subsection{Intégration avec le réentraînement}
\label{subsec:sprint3-cicd-retraining}

Le workflow GitHub Actions est complété par le mécanisme de
réentraînement via webhook. Lorsqu'un push modifie le code
d'entraînement ou les données, la plateforme détecte le changement
et relance automatiquement l'entraînement avec les nouvelles
versions.

Cette boucle ferme le cycle de vie : modification du code,
construction automatique, réentraînement, évaluation et déploiement
si le nouveau modèle est meilleur.


% =========================================================
% 5.10 BILAN DU SPRINT 3
% =========================================================

\section{Bilan du Sprint 3}
\label{sec:sprint3-bilan}

À la fin du Sprint 3, la plateforme disposait d'un cycle de vie
complet pour les modèles Machine Learning.

Le versionnement assurait la traçabilité de chaque modèle via le
MLflow Model Registry, avec une promotion automatique guidée par le
score global.

Le déploiement Blue/Green sur Kubernetes garantissait un service
continu sans interruption lors des mises à jour.

Les services de prédiction internes et externes permettaient
d'utiliser les modèles déployés depuis n'importe quelle application.

La gestion des accès contrôlait la visibilité des déploiements selon
le rôle et le partage.

Le monitoring surveillait en continu la santé des modèles et du
cluster, avec des alertes automatisées.

Enfin, les deux stratégies de réentraînement assuraient la
maintenance continue des modèles, que ce soit par réaction à une
dérive détectée ou par evolution du code source.


% =========================================================
% 5.11 CONCLUSION
% =========================================================

\section{Conclusion}
\label{sec:ch5-conclusion}

Ce troisième sprint a permis de boucler le cycle de vie complet des
modèles, depuis l'entraînement jusqu'au déploiement en production,
en passant par le versionnement, la prédiction et le monitoring.

La plateforme est désormais capable de gérer automatiquement
l'évolution des modèles dans le temps, grâce aux mécanismes de
réentraînement et d'alertes.

Le chapitre suivant sera consacré au déploiement sur Azure (AKS,
ACR, Key Vault) et à la mise en production finale de la plateforme.

% =========================================================
% FIN DU CHAPITRE 5
% =========================================================
